%================================================================
\chapter{Introduction}
\addcontentsline{toc}{chapter}{Introduction}
\markboth{Introduction}{Introduction}
%================================================================

\begin{infobox}
This document is \textbf{not} intended to be read from cover to cover.
Faculty and staff who are already comfortable with \LaTeX{}, Git, and
Python may proceed directly to the course-unit chapters. Those new to
one or more of the supporting tools should consult the relevant
infrastructure chapter first. Use the routing table on the following
pages to identify which chapters apply to your background.
\end{infobox}

%================================================================
\section*{Conventions Used in This Document}
%================================================================

\noindent
\textbf{Callout boxes.} Four types of callout box appear in the text:

\begin{itemize}
    \item \textbf{Information} (teal border): reference material and
          neutral context.
    \item \textbf{Tip} (blue border): practical guidance and
          shortcuts.
    \item \textbf{Warning} (magenta border): critical cautions
          that, if ignored, will cause errors or data loss.
    \item \textbf{Note} (orange border): items requiring attention
          but not critical.
\end{itemize}

\noindent
\textbf{Code listings.} All commands intended to be typed in a
terminal are shown in monospace code blocks. Windows commands use a
dark background; Linux/macOS commands use a lighter background.
Commands that work identically on all platforms use the neutral
listing style.

\noindent
\textbf{Cross-references.} Blue underlined text is a clickable
hyperlink. References to chapters, sections, figures, and tables are
also clickable in the PDF.

%================================================================
\section*{How This Document Is Organized}
%================================================================

The document is divided into two parts. The first part covers the
technical infrastructure required to work with the curriculum
repository: setting up the software environment, using the
command-line interface, writing in \LaTeX{}, managing versions with
Git and GitHub, and working with local language-model tools. These
chapters are written for faculty and staff who may not have a
background in software development; experienced contributors may skim
or skip them entirely.

The second part contains the course-unit chapters. Each unit
corresponds to one course in the department catalog and documents the
course title, student learning outcomes (SLOs), topical outline, and
assessment strategy. Units are numbered by their catalog identifier
(for example, unit 1023 corresponds to \textsc{Math 1023}). The
chapters within each part are independent; a reader interested only
in a specific course need not read any other unit chapter.

%================================================================
\section*{Chapter Dependencies}
%================================================================

      \renewcommand{\thefigure}{\arabic{figure}}
      \setcounter{figure}{0}
      Figure~\ref{fig:dependencies} shows the dependency structure of the
      infrastructure chapters. A solid arrow indicates a hard dependency:
      the target chapter requires the source to be completed first. A
      dashed arrow indicates recommended order only.

\begin{figure}[htbp]
\centering
\begin{tikzpicture}[
    node distance=1.6cm and 2.4cm,
    every node/.style={font=\small},
    req/.style={rectangle, rounded corners, draw=colorRequired,
                fill=colorRequired!15, text width=3.2cm, align=center,
                minimum height=0.9cm},
    ess/.style={rectangle, rounded corners, draw=colorEssential,
                fill=colorEssential!15, text width=3.2cm, align=center,
                minimum height=0.9cm},
    des/.style={rectangle, rounded corners, draw=colorDesirable,
                fill=colorDesirable!15, text width=3.2cm, align=center,
                minimum height=0.9cm},
    opt/.style={rectangle, rounded corners, draw=colorOptional,
                fill=colorOptional!15, text width=3.2cm, align=center,
                minimum height=0.9cm},
    hardarrow/.style={->, thick},
    softarrow/.style={->, dashed}
]
% Row 1
\node[req] (env)  {Environment Setup};
\node[req, right=of env] (cli) {CLI};
\node[req, right=of cli] (prog) {Intro to Programming};

% Row 2
\node[ess, below=of env] (pyenv) {Python Environments};
\node[des, below=of cli] (numops) {Numerical Operations};
\node[req, below=of prog] (latex) {\LaTeX{}};

% Row 3
\node[req, below=of pyenv] (git) {Git and GitHub};
\node[opt, below=of numops] (llm) {Local LLM};

% Arrows
\draw[hardarrow] (env) -- (cli);
\draw[hardarrow] (cli) -- (prog);
\draw[hardarrow] (env) -- (pyenv);
\draw[hardarrow] (prog) -- (pyenv);
\draw[softarrow] (pyenv) -- (numops);
\draw[hardarrow] (pyenv) -- (git);
\draw[softarrow] (cli) -- (latex);
\draw[softarrow] (numops) -- (llm);
\draw[softarrow] (git) -- (llm);
\end{tikzpicture}
\caption{Dependency graph for the infrastructure chapters. Solid arrows
indicate hard dependencies; dashed arrows indicate recommended order. Node
color encodes category: \textcolor{colorRequired}{\textbf{Required}},
\textcolor{colorEssential}{\textbf{Essential}},
\textcolor{colorDesirable}{\textbf{Desirable}},
\textcolor{colorOptional}{\textbf{Optional}}.}
\label{fig:dependencies}
\end{figure}


\begin{landscape}
%================================================================
\section*{Which Chapters Do You Need?}
%================================================================

Use Table~\ref{tab:routing} to identify your reading path based on your
familiarity with the supporting tools. Chapters marked \textbf{Required}
must be completed before working with the repository. \textbf{Skim} means
read the section headings and run the exercises but skip the explanatory
prose. \textbf{Skip} means the chapter is optional for your profile.

\begin{table}[htbp]
\centering
\caption{Recommended reading path by prior experience with the toolchain.}
\label{tab:routing}
\footnotesize
\begin{tabular}{l*{8}{c}}
\toprule
\textbf{Profile}
  & \rotatebox{70}{Env.\ Setup}
  & \rotatebox{70}{CLI}
  & \rotatebox{70}{Intro Prog.}
  & \rotatebox{70}{Python Envs.}
  & \rotatebox{70}{Numerical Ops.}
  & \rotatebox{70}{\LaTeX{}}
  & \rotatebox{70}{Git}
  & \rotatebox{70}{LLM} \\
\midrule
New to all tools
  & $\bullet$ & $\bullet$ & $\bullet$ & $\bullet$ & $\bullet$
  & $\bullet$ & $\bullet$ & S \\
Familiar with \LaTeX{}, new to Python and Git
  & $\bullet$ & $\bullet$ & $\bullet$ & $\bullet$ & S
  & -- & $\bullet$ & S \\
Familiar with Python and Git, new to \LaTeX{}
  & $\bullet$ & S & S & S & --
  & $\bullet$ & -- & S \\
Experienced with all tools
  & S & -- & -- & S & --
  & -- & -- & S \\
\bottomrule
\end{tabular}

\bigskip
\noindent\footnotesize
$\bullet$ = Required \quad S = Skim \quad -- = Skip
\normalsize
\end{table}

%================================================================
\section*{Progress Record}
%================================================================

The final chapter, Chapter~\ref{ch:assessment} (Assessment and Progress
Record), organizes the key completion items from all infrastructure chapters
into four categories: Required, Essential, Desirable, and Optional. Open
the fillable PDF form in Adobe Acrobat Reader to record and save your
responses. Use it throughout onboarding as a readiness checklist.

\newpage
%================================================================
\section*{Estimated Completion Times}
%================================================================

Table~\ref{tab:times} shows estimated completion times per infrastructure
chapter. Times assume no prior exposure to the topic; contributors with
background in a given area will complete it faster. Course-unit chapters
are read as reference material and are not included in this table; reading
time for a single unit is approximately 20 to 30 minutes.

\begin{table}[htbp]
\centering
\caption{Estimated completion time per infrastructure chapter.}
\label{tab:times}
\begin{tabular}{clr}
\toprule
\textbf{Chapter} & \textbf{Title} & \textbf{Estimated Time} \\
\midrule
1  & Environment Setup                              & 2 hours \\
2  & Command-Line Interfaces and Shell Scripting    & 2 hours \\
3  & Introduction to Programming Environments       & 2 hours \\
4  & Python Virtual Environments and Workflow Automation & 1.5 hours \\
5  & Numerical Operations with NumPy                & 2 hours \\
6  & Introduction to \LaTeX{}                       & 1.5 hours \\
7  & Version Control with Git and GitHub            & 1.5 hours \\
8  & Local LLM Infrastructure with Ollama           & 1 hour (contributors) / 3 hours (administrators) \\
9  & Assessment and Progress Record                 & Ongoing \\
\midrule
   & \textbf{Total (full path)}                     & \textbf{$\approx$ 13.5--15.5 hours} \\
\bottomrule
\end{tabular}
\end{table}

\end{landscape}